
\subsection*{Abstract}
We propose \textbf{Community-Structured Attention (CSA)}, an attention mechanism that enables multi-hop reasoning over large Knowledge Graphs (KGs) without the $O(N^2)$ complexity of global attention matrices. CSA maps the structural components of the Transformer architecture \cite{vaswani2017attention} directly onto graph operations, utilizing community partitions as discrete "Attention Heads." We define a unified scoring function that integrates semantic similarity, community-level topology, and structural centrality. Benchmark results on the \textbf{Hetionet} \cite{hetionet2017} biomedical dataset demonstrate that CSA achieves a Mean Reciprocal Rank (MRR) of \textbf{0.68}, a \textbf{+183\% improvement} over breadth-first search baselines. Furthermore, on the \textbf{MetaQA 3-hop} \cite{metaqa2017} reasoning task, CSA improves MRR by \textbf{+350\%}, demonstrating superior beam steering in deep multi-hop traversals while maintaining full "Glass-Box" interpretability. As of v2.51.0, the CSA formula has been expanded to 10 parameters covering temporal decay, node recency, synthesis-density penalty, and grounding confidence, with both batch (CSAParameterLearner) and online per-community (MetaParameterLearner) learning, achieving MetaQA canonical results of H@1=46.1\%/30.0\%/12.5\% across 1-, 2-, and 3-hop tasks.

\subsection*{1. Introduction}
The dominance of Transformer architectures in Natural Language Processing has inspired attempts to apply similar attention-based principles to graph structures. However, Graph Attention Networks (GATs) \cite{velickovic2018gat} typically operate on local ego-networks and struggle with global structural context. CSA addresses this by introducing a "Soft Community Constraint," where attention weights are influenced by the membership of nodes in pre-computed structural partitions (DSCF/TSC).

\subsection*{2. The Cerebrum Mapping}
CSA is built on a direct functional analogy to the Transformer \cite{vaswani2017attention}:
\begin{itemize}
  \item \textbf{Communities} act as \textbf{Attention Heads}, focusing the search on specific semantic neighborhoods.
  \item \textbf{Centrality Features} (PageRank, Betweenness) serve as \textbf{Positional Encodings}, providing structural context.
  \item \textbf{Traversal Paths} function as a \textbf{KV Cache}, memoizing the reasoning history.
\end{itemize}


\subsection*{3. Methodology}

\subsubsection*{3.1 The CSA Formula}
The attention weight $a(u,v,k)$ for an edge from $u$ to $v$ at hop $k$ is defined as:
$$\begin{aligned}
a(u,v,k) = \sigma( & \alpha \cdot \mathcal{S}_{sem}(u,v) + \beta \cdot \mathcal{S}_{com}(u,v) + \\
& \gamma \cdot w_{rel} - \delta \cdot d_{norm}(u,v) + \epsilon \cdot \phi(k) )
\end{aligned}$$

\subsubsection*{3.2 The Community Signal ($\mathcal{S}_{com}$)}
Unlike GATs \cite{velickovic2018gat} which treat all neighbors equally, CSA scales weights based on community topology:
\begin{itemize}
  \item \textbf{Intra-community}: $1.0$
  \item \textbf{Adjacent-community}: $0.5$
  \item \textbf{Distant-community}: $e^{-\lambda d_{com}}$
\end{itemize}


\subsection*{4. Enterprise Hardening (v2.51.0)}
The v2.51.0 release introduces \textbf{Adaptive Parameter Learning}, utilizing a \textbf{MetaParameterLearner} to autonomously adjust the $(\alpha, \beta, \gamma, \delta, \epsilon)$ coefficients per-community based on query feedback. This closes the gap between zero-shot and supervised performance without the need for global retraining.

\subsection*{5. Conclusion}
CSA provides a scalable, Interpretable AI (XAI) alternative to black-box graph embeddings. By grounding attention in the structural consensus of the graph, it enables complex multi-hop reasoning that is both computationally efficient and mathematically verifiable. In CEREBRUM v2.51.0, the 10-parameter CSA formula with online per-community learning achieves MetaQA H@1 of 46.1\% (1-hop), 30.0\% (2-hop), and 12.5\% (3-hop), alongside WebQSP H@1=6.27\%, H@10=20.84\%, and MRR=10.66\% --- establishing CSA as a competitive and interpretable alternative to embedding-based KG reasoning.

\medskip\noindent\rule{\linewidth}{0.4pt}\medskip

\section*{6. Recent Advances (v2.24.0 -> v2.51.0)}

The CSA formula and its associated learning infrastructure have undergone significant expansion since v2.24.0. The following describes the key advances relevant to this paper.

\textbf{10-Parameter CSA Formula (Phase 43/45).} The original 5-parameter formula \texttt{(alpha, beta, gamma, delta, epsilon)} covering semantic similarity, community score, edge-type weight, distance penalty, and hop decay has been extended to 10 parameters:

\begin{lstlisting}
a(u,v,k) = sigmoid(
    alpha   * sim          # semantic similarity (cosine)
  + beta    * cs           # community score (structural membership)
  + gamma   * etw          # edge-type weight
  - delta   * nd           # normalised distance penalty
  + epsilon * hd           # hop decay
  + zeta    * pr_v         # PageRank prior
  + eta     * td           # temporal decay
  + iota    * nr_v         # node recency
  - mu      * sd           # synthesis-density penalty
  + theta   * grounding    # confidence / grounding score
)
\end{lstlisting}

Default weights: \texttt{(0.4, 0.4, 0.1, 0.05, 0.05, 0.1, 0.1, 0.05, 0.1, 1.0)}. The synthesis-density penalty (\texttt{-mu*sd}) is particularly significant: it prevents over-reliance on REM-synthesized Synaptic Bridge edges, maintaining reasoning transparency.

\textbf{Unified ReasoningLogit Vector.} All 10 features are bundled into a \texttt{ReasoningLogit} dataclass that threads through scoring, learning, and feedback logging. This ensures that the full feature vector is available for inspection at every hop, preserving the "Glass-Box" property of the original design.

\textbf{Batch and Online Parameter Learning (Phase 22/45/48).} Two learning regimes are now fully implemented:
\begin{itemize}
  \item \texttt{CSAParameterLearner.fit()}: batch gradient descent over accumulated (positive, negative) path pairs, triggered via \texttt{POST /retrain}. Updates the global 10-parameter prior.
  \item \texttt{MetaParameterLearner}: online SGD per community, updated on every \texttt{POST /feedback} call. Enables community-specific adaptation without global retraining.
\end{itemize}


\textbf{Params Persistence (Phase 47).} \texttt{MetaParameterLearner.to\_dict()} / \texttt{from\_dict()} enables checkpoint and restore via \texttt{POST /params}. The \texttt{--params-file} CLI flag loads a checkpoint at server startup, enabling zero-downtime parameter rollback.

\textbf{Benchmark Results.}

\begin{table}[h]
\small\centering
\begin{tabular}{lll}
\toprule
Dataset & Metric & v2.51.0 \\
\midrule
MetaQA 1-hop & H@1 / H@10 & 46.1% / 96.6% \\
MetaQA 2-hop & H@1 / H@10 & 30.0% / 86.3% \\
MetaQA 3-hop & H@1 / H@10 & 12.5% / 50.3% \\
WebQSP OPT & H@1 / H@10 / MRR & 6.27% / 20.84% / 10.66% \\
\bottomrule
\end{tabular}
\end{table}


\section*{7. Phase 55 Advances}

\textbf{GraphSAGE Neighbourhood Smoothing (Phase 55).} \texttt{smooth\_with\_graphsage(embeddings, G)} applies a one-pass mean neighbourhood aggregation after base encoding:

$$\tilde{\mathbf{e}}_v = \frac{1}{1+|\mathcal{N}(v)|}\left(\mathbf{e}_v + \sum_{u \in \mathcal{N}(v)} \mathbf{e}_u\right)$$

The enriched embeddings make the \texttt{alpha} (semantic similarity) term in the CSA formula significantly more discriminating --- nodes in the same community share more similar neighbourhood-aggregated representations. \texttt{CerebrumGraph.build(use\_graphsage=True)} enables smoothing automatically after base encoding. Complexity is $O(|E| \times d)$ where $d$ is the embedding dimension.

\textbf{TemporalCalibrator (Phase 55).} Grid-searches \texttt{eta} (temporal decay) and \texttt{iota} (node recency) against a labelled validation set to maximise Recall@K. The \texttt{calibrate()} method enumerates a parameter grid, calls \texttt{measure\_recall()} at each point, and applies the best-found parameters to the CSAEngine. A \texttt{try/finally} block guarantees original parameters are restored if calibration is interrupted --- ensuring that a failed calibration run never leaves the CSAEngine in a partially-modified state.

\textbf{Engram-Steered Traversal (Phase 55).} \texttt{Engram} tracks relation-sequence patterns from previous successful Engram traces. \texttt{EngramTraversal.\_prune\_candidates()} applies:

$$s_\text{eff}(c) = s(c) \times (1 + \lambda_\text{engram} \cdot \text{affinity}(\text{rel\_seq}))$$

where \texttt{affinity} is derived from accumulated \texttt{\_counts}. This biases beam search toward known-productive reasoning chains without modifying graph structure. The cache is durable --- \texttt{save(path)} serializes to JSON and \texttt{load(path)} restores counts on restart, so learned relation patterns survive process restarts.